% ####################################################################
% Basic configuration and packages
% ------------------------------------------------------------------------------------------------------------------
% Package for discovering wrong and outdated usage of LaTeX.
% More information to be found in l2tabu English version.
\RequirePackage[l2tabu, orthodox]{nag}
% Class of LaTeX document: {size of paper, size of font}[document class]
\documentclass[a4paper,11pt]{article}



% ------------------------------------------------------
% Packages not tied to particular normal language
% ------------------------------------------------------
% This package should improved spaces in the text.
\usepackage{microtype}
% Add few important symbols, like text Celcius degree
\usepackage{textcomp}



% ------------------------------------------------------
% Polonization of LaTeX document
% ------------------------------------------------------
% Basic polonization of the text
\usepackage[MeX]{polski}
% Switching on UTF-8 encoding
\usepackage[utf8]{inputenc}
% Adding font Latin Modern
\usepackage{lmodern}
% Package is need for fonts Latin Modern
\usepackage[T1]{fontenc}



% ------------------------------------------------------
% Setting margins
% ------------------------------------------------------
\usepackage[a4paper, total={14cm, 25cm}]{geometry}



% ------------------------------------------------------
% Setting vertical spaces in the text
% ------------------------------------------------------
% Setting space between lines
\renewcommand{\baselinestretch}{1.1}

% Setting space between lines in tables
% \renewcommand{\arraystretch}{1.4}



% ------------------------------------------------------
% Packages for scientific papers
% ------------------------------------------------------
% Switching off \lll symbol, that I guess is representing letter "Ł"
% It collide with `amsmath' package's command with the same name
% \let\lll\undefined
% Basic package from American Mathematical Society (AMS)
% \usepackage[intlimits]{amsmath}
% Equations are numbered separately in every section.
% \numberwithin{equation}{section}

% Other very useful packages from AMS
% \usepackage{amsfonts}
% \usepackage{amssymb}
% \usepackage{amscd}
% \usepackage{amsthm}

% Package for writting physical units
% \usepackage{siunitx}

% Package with better looking calligraphy fonts
% \usepackage{calrsfs}

% Package with better looking greek letters
% Example of use: pi -> \uppi
% \usepackage{upgreek}
% Improving look of lambda letter
% \let\oldlambda\Lambda
% \renewcommand{\lambda}{\uplambda}




% ------------------------------------------------------
% BibLaTeX
% ------------------------------------------------------
% Package biblatex, with biber as its backend, allow us to handle
% bibliography entries that use Unicode symbols outside ASCII.
\usepackage[
language=polish,
backend=biber,
style=alphabetic,
url=false,
eprint=true,
]{biblatex}

\addbibresource{Słownik-informatyczno-polski-Bibliography.bib}





% ------------------------------------------------------
% Defining new environments (?)
% ------------------------------------------------------
% Defining enviroment "Wniosek"
\newtheorem{corollary}{Wniosek}
\newtheorem{definition}{Definicja}
\newtheorem{theorem}{Twierdzenie}





% ------------------------------------------------------
% Local packages
% You need to put them in the same directory as .tex file
% ------------------------------------------------------
% Package containing various command useful for working with a text
% \usepackage{./Local-packages/general-commands}





% ------------------------------------------------------
% Package "hyperref"
% They advised to put it on the end of preambule
% ------------------------------------------------------
% It allows you to use hyperlinks in the text
\usepackage{hyperref}










% ####################################################################
% Title and author of the text
\title{Słownik systemy operacyjne~-- język polski}

% \author{}


% \date{}
% ####################################################################










% ####################################################################
% Beginning of the document
\begin{document}
% ####################################################################





% ######################################
\maketitle
% ######################################





Przy tworzeniu tego słonika korzystaliśmy z~następujących materiałów. %
% Bruce Eckel \textit{Thinking in~Java. Edycja
%   polska, wydanie~III},
% \parencite{Eckel-Thinking-in-Java-Ed-polska-Wyd-III-Pub-2003}.

\vspace{1.5em}



\noindent
% \textbf{Shell}~-- modyfikatory ????dostępu. \\
\textbf{Busy waiting}~-- aktywne oczekiwanie. \\
\textbf{Cache hit}~-- trafienie pamięci podręcznej. \\
\textbf{Cache line}~-- linia pamięci podręcznej. \\
\textbf{Core dump}~-- zapis pamięci, zrzut pamięci rdzenia. \\
\textbf{CPU (Central Processing Unit)}~-- procesor, centralna jednostka
obliczeniowa. \\
\textbf{DMA (Direct Memory Access)}~-- bezpośredni dostęp do pamięci. \\
\textbf{GUI (Graphical User Interface)}~-- graficzny interfejs użytkownika. \\
\textbf{Heap}~-- sterta. \\
\textbf{Hypervisor}~-- hipernadzorca. \\
% \textbf{Concurrency}~-- współbieżność. \\
\textbf{IDE (Integrated Drive Electronics)}~-- . \\
\textbf{Kernel mode}~-- tryb jądra. \\
\textbf{MMU (Memory Menagement Unit)}~-- moduł zarządzania pamięcią. \\
\textbf{Motherboard}~-- płyta główna, płyta macierzysta. \\
\textbf{LSI (Large Scale Integration)}~-- integracja wielkiej skali. \\
\textbf{Line printer}~-- drukarka wierszowa. \\
\textbf{PSW (Program Status Word)}~-- słowo stanu programu. \\
\textbf{Pipe}~-- potok. \\
\textbf{RAM (Random Access Memory)}~-- pamięć o~stały czasie dostępu do
losowych elementów. \\
\textbf{ROM (Read Only Memory)}~-- pamięć tylko do odczytu. \\
\textbf{Shell}~-- powłoka. \\
\textbf{Spooling (simultaneous peripheral operatio on line)}~-- jednoczesne
działanie podłączonych urządzeń. \\
\textbf{Stack}~-- stos. \\
\textbf{Supervisor mode}~-- tryb nadzorcy. \\
\textbf{Timesharing}~-- podział czasu. \\
\textbf{User mode}~-- tryb użytkownika. \\



% \textbf{Package access}~-- dostęp pakietowy. \\
% \textbf{Parametrized type}~-- typ parametryzowalny. \\
% \textbf{Persistent}~-- trwałość. \\
% \textbf{Substitution principle}~-- zasada zastępowalności. \\
% \textbf{Throwing an exception}~-- rzucanie wyjątków. \\
% \textbf{Translation unit} (alt.~\textbf{compilation unit})~-- jednostka
% tłumaczenia/translacji. \\
% \textbf{Upcasting}~-- rzutowanie w~górę. \\
% Klasa potomna, \\

% \textbf{Generics}~-- generyki. \\
% \textbf{Implementation hiding} (alt.~\textit{encapsulation})~--
% ukrywanie implementacji. \\










% ######################################
\section{Nazwy konkretnych programów i~powiązane rzeczy}
% ######################################


\textbf{CP/M (Control Program for Microcomputers)}~-- . \\
\textbf{CTSS (Compatible Time Sharing System)}~-- . \\
\textbf{DOS (Disc Operation System)}~-- . \\
\textbf{JCL (Job Control Language)}~-- . \\
\textbf{MAD (Michigan Algorithm Decoder)}~-- . \\
\textbf{MULTICS (MULTiplexed Information and Computing Service)}~--. \\
\textbf{MS-DOS (MicroSoft Disk Operating System)}~-- . \\
\textbf{Windows NT (Windows New Technology)}~-- . \\
\textbf{Windows Me (Windows Millenium)}~-- . \\
% \textbf{Parametrized type}~-- typ parametryzowalny. \\
% \textbf{Persistent}~-- trwałość. \\
% \textbf{Substitution principle}~-- zasada zastępowalności. \\
% \textbf{Throwing an exception}~-- rzucanie wyjątków. \\
% \textbf{Translation unit} (alt.~\textbf{compilation unit})~-- jednostka
% tłumaczenia/translacji. \\
% \textbf{Upcasting}~-- rzutowanie w~górę. \\
% Klasa potomna, \\







% ####################################################################
% ####################################################################
% Bibliography

\printbibliography





% ############################
% End of the document

\end{document}
